\documentclass[11pt]{article}
\usepackage[T1]{fontenc}
\usepackage{lmodern,amsmath,amssymb,amsthm}
\usepackage[margin=1in]{geometry}
\usepackage[hidelinks]{hyperref}
\setlength{\emergencystretch}{3em}
\newtheorem{proposition}{Proposition}
\title{Translated periods and integral Hecke congruences at conductor 234446}
\author{\small\href{mailto:admin@cosmosgeek.org}{admin@cosmosgeek.org}}
\date{Working draft for specialist review}
\begin{document}
\maketitle
\begin{abstract}
For the rank-four elliptic curve of conductor 234446, we study two
ways of recovering its second complex central derivative. On adjoining
the two rank-three newforms at the same conductor, the Fourier-integral
lattice contains the normalized eigenform lattice with index four. We
compute its Hecke action and the associated genus-two gluings. A
Dirichlet-series transfer converges on compact sets but fails an
absolute second-moment bound.
A translated logarithmic period \(P\) satisfies \(2\pi P(0)=L''(E,1)\).
All real translates \(P(s)\) for \(s\ge1\) are positive, while their
inverse Laplace density changes sign infinitely often. We give a signed
reconstruction of the endpoint and compute a closed modular-symbol
period for the quadratic character of conductor three. The central order
on the rank-four curve remains undetermined by these identities.
\end{abstract}

\section{A fixed-level transfer at conductor 234446}
\label{sec:234446-transfer}

\subsection{Models and analytic hypotheses}

Put \(N=234446=2\cdot117223\), and consider the elliptic curves
\[
\begin{aligned}
E &: [1,-1,0,-79,289],\\
A &: [1,1,0,-696,6784],\\
C &: [1,1,1,-949,-7845].
\end{aligned}
\]
Cremona's current curve and optimality tables \cite{Cremona} list these
models as the sole members of isogeny classes \(a,b,c\) at this level,
with Mordell--Weil ranks \(4,3,3\), trivial rational torsion, and Manin
constant one. We distinguish them by their equations.
Let \(f_E,f_A,f_C\in S_2(\Gamma_0(N))\) denote the normalized
rational newforms. For \(p=3,11\), count points using the quadratic character of
\(4x^3+b_2x^2+2b_4x+b_6\). At \(117223\) use the Legendre symbols of
\(-c_6\), and at \(2\) count directly on each reduction. This gives
\[
\begin{array}{c|rrrr}
 &a_2&a_3&a_{11}&a_{117223}\\ \hline
E&-1&-3&-6& 1\\
A&-1&-1&-4&-1\\
C& 1&-1&-6& 1
\end{array}
\]
and hence global root numbers \(+1,-1,-1\), respectively.

Write
\[
\Lambda_F(s)=
 \left(\frac{\sqrt N}{2\pi}\right)^s
 \Gamma(s)L(F,s),\qquad
g_F(t)=e^t f_F(ie^t/\sqrt N).
\]
The functional equations imply that \(g_E\) is even and
\(g_A,g_C\) are odd. With the standard weight-two Fricke
normalization,
\[
W_Nf_E=-f_E,\qquad W_Nf_A=f_A,\qquad W_Nf_C=f_C.
\]
By the analytic-rank-zero-and-one theorems of
Gross--Zagier--Kolyvagin \cite{GZK}, the established algebraic ranks imply
\[
\operatorname{ord}_{s=1}L(E,s)\ge2,\qquad
\operatorname{ord}_{s=1}L(A,s),
\operatorname{ord}_{s=1}L(C,s)\ge3.
\]
For \(M_j(g)=\int_{\mathbb R}t^jg(t)\,dt\), Mellin transformation gives
\(M_j(g_F)=\Lambda_F^{(j)}(1)\). Consequently
\[
M_2(g_A')=M_2(g_C')=0,\qquad
M_2(g_A/t)=M_2(g_C/t)=0.
\]
The congruence and the correspondences below test transfer of these moment identities to the even kernel of \(E\).

\subsection{Exact cubic congruences}

For a Weierstrass model
\(y^2+a_1xy+a_3y=x^3+a_2x^2+a_4x+a_6\),
put \(u=4x\). Its nonidentity \(2\)-torsion points are
parametrized by the roots of
\[
F(u)=u^3+b_2u^2+8b_4u+16b_6,
\]
where
\(b_2=a_1^2+4a_2\),
\(b_4=a_1a_3+2a_4\), and
\(b_6=a_3^2+4a_6\).

For the three models this yields
\[
\begin{aligned}
F_E(u)&=u^3-3u^2-1264u+18496,\\
F_A(u)&=u^3+5u^2-11136u+434176,\\
F_C(u)&=u^3+5u^2-15176u-502064.
\end{aligned}
\]

\begin{proposition}
There are Galois-equivariant isomorphisms
\(E[2]\simeq A[2]\simeq C[2]\), explicitly induced on
the nonzero \(2\)-torsion coordinates by
\[
\boxed{
\begin{aligned}
F_A(864-20u-u^2)&\equiv0\pmod{F_E(u)},\\
F_C(860-16u-u^2)&\equiv0\pmod{F_E(u)}.
\end{aligned}}
\]
\end{proposition}

\begin{proof}
The two identities follow by polynomial division over
\(\mathbb Z[u]\). Irreducibility and separability can be checked
without a discriminant calculation. Modulo \(3\), the cubics are
\[
\begin{aligned}
\overline F_E&=u^3+2u+1,\\
\overline F_A&=u^3+2u^2+1,\\
\overline F_C&=u^3+2u^2+u+1.
\end{aligned}
\]
Their value triples at \(u=0,1,2\) are, respectively,
\[
(1,1,1),\qquad(1,1,2),\qquad(1,2,1).
\]
None has a root in \(\mathbb F_3\), so all three cubics
are irreducible over \(\mathbb Q\). In characteristic zero
they therefore have three distinct roots.

Let \(\alpha\) be a root of \(F_E\).
The displayed substitutions produce roots of the two other
irreducible cubics in \(\mathbb Q(\alpha)\).
Each resulting root has degree three over \(\mathbb Q\);
hence the substitutions induce isomorphisms of cubic fields.
Applying the three embeddings shows that they biject
the respective root sets, with no collapsed roots, and
commute with Galois. Including the identity point gives
the asserted isomorphisms of \(2\)-torsion Galois modules.
\end{proof}

At good primes these isomorphisms give congruent Frobenius
traces modulo \(2\). At the two multiplicative primes all
three local eigenvalues are \(\pm1\), so their reductions
also agree. The Euler recurrences and multiplicativity
therefore yield
\[
f_E\equiv f_A\equiv f_C\pmod2
\]
coefficientwise, including the bad Euler factors. The traces at 3 already differ modulo 4, so
the displayed identifications of two-torsion do not lift to
identifications of four-torsion.

\subsection{The saturated Fourier lattice and its Hecke action}

Define integral cusp forms
\[
h_A=\frac{f_E-f_A}{2},\qquad
h_C=\frac{f_E-f_C}{2},
\]
and let
\[
\mathcal V=\mathbb Q f_E+\mathbb Q f_A+\mathbb Q f_C,\qquad
\mathcal L=\mathbb Zf_A+\mathbb Zh_A+\mathbb Zh_C.
\]

\begin{proposition}
The lattice \(\mathcal L\) is precisely
\[
\mathcal V\cap\mathbb Z[[q]],
\]
and the sublattice generated by the three normalized
eigenforms has index \(4\) in \(\mathcal L\).
\end{proposition}

\begin{proof}
All three displayed generators have integral Fourier
coefficients. Their \(q^1,q^3,q^{11}\) coefficient matrix is
\[
\begin{pmatrix}
1&0&0\\
-1&-1&-1\\
-4&-1&0
\end{pmatrix},
\qquad \det=-1.
\]
Thus any rational combination with integral Fourier
coefficients has integral coordinates in this basis,
proving saturation.

Since
\[
f_E=f_A+2h_A,\qquad
f_C=f_A+2h_A-2h_C,
\]
the change-of-basis determinant between the normalized
eigenform lattice and \(\mathcal L\) has absolute value \(4\).
\end{proof}

In the ordered basis \((f_A,h_A,h_C)\), the exact matrices
are
\[
T_3=
\begin{pmatrix}
-1&-1&-1\\
0&-3&-2\\
0&0&-1
\end{pmatrix},\qquad
T_{11}=
\begin{pmatrix}
-4&-1&0\\
0&-6&0\\
0&0&-6
\end{pmatrix},
\]
and
\[
W_N=
\begin{pmatrix}
1&-1&-1\\
0&-1&-2\\
0&0&1
\end{pmatrix}.
\]
These follow immediately from the eigenvalues in the
preceding coefficient table and the definitions of
\(h_A,h_C\). In particular,
\[
\overline{T_3}=\overline{W_N}=
\begin{pmatrix}
1&1&1\\
0&1&0\\
0&0&1
\end{pmatrix}
\quad\text{over }\mathbb F_2.
\]
The integral reduction is therefore not a direct sum
of three independently split eigenlines.

Over \(\mathbb Q\), however, \(T_3\) has eigenvalue
\(-3\) on \(f_E\) and \(-1\) on both \(f_A,f_C\).
The spectral projector onto the \(E\)-eigenline is
\[
\boxed{e_E=-\frac{T_3+1}{2}.}
\]
It is not integral on \(\mathcal L\), since
\[
e_E(h_A)=e_E(h_C)=\frac{f_E}{2}
\notin\mathcal L.
\]
The exact integral identities
\[
(T_3+1)h_A=(T_3+1)h_C=-f_E
\]
therefore transfer the mixed congruence vectors to
\(f_E\), but do not transfer \(f_A\) or \(f_C\):
\[
(T_3+1)f_A=(T_3+1)f_C=0.
\]
The mixed vectors already contain the unknown
\(f_E\)-component. This is the first failed step in a
characteristic-zero transfer of the annihilated
rank-three moments. Any characteristic-zero Hecke-equivariant homomorphism from either of the other eigenlines to the \(E\)-line vanishes, by the \(T_3\)-eigenvalue.

\subsection{Genus-two Jacobians from gluing}

Let \(F\) denote either \(A\) or \(C\), and let
\(\phi_F:E[2]\to F[2]\) be the corresponding explicit
Galois-module isomorphism. Define
\[
\Gamma_F=\{(P,\phi_F(P)):P\in E[2]\},
\qquad
J_F=(E\times F)/\Gamma_F,
\]
with quotient isogeny \(\pi_F:E\times F\to J_F\).

Every isomorphism between the two symplectic
\(\mathbb F_2\)-spaces is anti-symplectic at level \(2\).
Thus \(\Gamma_F\) is maximal isotropic.

The curves \(E\) and \(F\) are not geometrically
isogenous. Indeed, at the good prime \(3\),
\(a_3(E)=-3\) gives supersingular reduction, whereas
\(a_3(F)=-1\) gives ordinary reduction. These reduction
types are invariant under geometric isogeny.

Howe--Lepr\'evost--Poonen, Proposition 3, therefore
applies: each principally polarized \(J_F\) is the
Jacobian of a smooth genus-two curve over \(\mathbb Q\).
See \cite[Proposition~3]{HLP}.

Write \(\lambda_E,\lambda_F\) for the canonical elliptic
principal polarizations and \(\lambda_{J_F}\) for the
descended principal polarization. Then
\[
\boxed{
\deg\pi_F=4,\qquad
\pi_F^*\lambda_{J_F}
=2(\lambda_E\oplus\lambda_F).
}
\]
Equivalently, the polarization adjoint satisfies
\(\pi_F^\dagger\pi_F=[2]\).

Consequently,
\[
\operatorname{rank}J_F(\mathbb Q)=4+3=7.
\]
Let
\[
I_F=
\left[
J_F(\mathbb Q)/J_F(\mathbb Q)_{\rm tors}:
\pi_F((E\times F)(\mathbb Q))
\ \bmod J_F(\mathbb Q)_{\rm tors}
\right].
\]
This index is finite. The adjoint identity shows that
its cokernel is killed by \(2\), so \(I_F\) is a power
of \(2\). By functoriality of canonical heights, the
height pairing on the image of the product point
lattice is twice the product pairing. Hence
\[
\boxed{
\operatorname{Reg}(J_F)=
\frac{2^7\operatorname{Reg}(E)
                 \operatorname{Reg}(F)}
     {I_F^2}.
}
\]
This holds separately for \(J_A\) and \(J_C\), with
their actual respective Mordell--Weil indices.

At the infinite place the isogeny similarly enlarges
the product period lattice by the graph of the
\(2\)-torsion identification:
\[
\Lambda_E\oplus\Lambda_F
\subset\Lambda_{J_F}
\subset\tfrac12(\Lambda_E\oplus\Lambda_F),
\qquad
[\Lambda_{J_F}:\Lambda_E\oplus\Lambda_F]=4,
\]
under compatible complex uniformizations.
Thus integral half-periods are genuinely mixed,
while rational cohomology still splits:
\[
H^1(J_F,\mathbb Q)
\simeq H^1(E,\mathbb Q)\oplus H^1(F,\mathbb Q).
\]
Since \(\operatorname{Hom}_{\overline{\mathbb Q}}(E,F)=0\),
algebraic correspondences cannot provide a nonzero
off-diagonal transfer between the elliptic factors.

\subsection{Consequence for the second derivative}

Isogeny invariance gives
\[
L(J_F,s)=L(E,s)L(F,s).
\]
For the corresponding completed functions, the
established lower bounds on vanishing imply the exact
identity
\[
\boxed{
\Lambda_{J_F}^{(5)}(1)
=10\,\Lambda_E''(1)\Lambda_F'''(1).
}
\]
If \(\Lambda_F'''(1)\ne0\), this identifies the
vanishing of the original second derivative with
fifth-derivative vanishing on \(J_F\). That
nonvanishing condition is additional to the proven
algebraic rank-three statement.

The regulator formula, integral Hecke extension and
glued period lattices do not themselves force the
fifth derivative to vanish. The fifth derivative of the glued Jacobian remains undetermined by this calculation.

\section{Transfer by the ratio of two Dirichlet series}
For the pair \(E,A\) write
\[
L(E,s)/L(A,s)=\sum_{m\ge1} b_m m^{-s}\qquad(\Re s>3/2).
\]
Both Dirichlet series and the reciprocal Euler product converge
absolutely there. Put \(h_m(t)=g_A(t+\log m)-g_A(t-\log m)\).
Then
\begin{equation}\label{eq:ratio}
g_E(t)=\sum_{m\ge1}\frac{b_m}{2m}h_m(t),\qquad
\int_{\mathbb R}t^2h_m(t)\,dt=-4\log(m)\Lambda_A'(1)=0 .
\end{equation}
Indeed \(a_n(E)=\sum_{m\mid n}b_ma_{n/m}(A)\).
Substitution into the exponentially convergent newform expansion gives
\(g_E(t)=\sum_m(b_m/m)g_A(t+\log m)\). Average this with its reflection
at \(-t\), using the opposite parities of the kernels. The series
converge absolutely and uniformly for \(t\) in a compact interval:
the translated terms decay exponentially in \(m\) there. Translation
in the moment integral gives the second identity.

\begin{proposition}
For this pair,
\[
\sum_m\frac{|b_m|}{2m}\int_{\mathbb R}t^2|h_m(t)|\,dt=\infty .
\]
\end{proposition}
\begin{proof}
As \(a=\log m\to\infty\), the integral divided by \(a^2\) tends to
\(2\int|g_A|\). On each of the two windows \(|t\mp a|\le R\), the
other translate tends uniformly to zero and \(t^2/a^2\to1\).
Letting \(R\to\infty\) gives the lower bound. The triangle inequality
and rapid decay give the matching upper bound.
At a good prime \(p\), \(b_p=a_p(E)-a_p(A)\), and point counts give
\(a_5(E)=-4,\ a_5(A)=-2\). The Frobenius class of \(5\) in the
compositum of the two \(3\)-division fields consequently supplies, by
Chebotarev, a positive-density set of primes with \(b_p\not\equiv0
\bmod3\). Partial summation shows that the sum of
\((\log p)^2/p\) over that set diverges.
\end{proof}
The second moments of finite partial sums in \eqref{eq:ratio} all
vanish. Compact convergence supplies \(\Lambda_E''(1)\) if one first
takes the limit of integrals over bounded intervals and only afterwards
enlarges the interval. Termwise integration would require additional
signed tail control.

For comparison, if \(L(A,s)\) has exact analytic order three, the
completed functional equations give an odd quotient and
\[
\operatorname*{Res}_{u=0}\frac{L(E,1+u)}{L(A,1+u)}
=\frac{3L''(E,1)}{L'''(A,1)}.
\]
At \(q=117223\) the bad-prime quotient is \((1+q^{-s})/(1-q^{-s})\),
a unit at the center. The Atkin--Lehner projector \((1-W_q)/2\)
separates \(f_E\) from \(f_A\) with Petersson norm one. It gives no
bound for the reciprocal of the displayed residue.

\section{The translated logarithmic period}
Write \(F(y)=f_E(iy)\), and use the principal logarithm in
\[
P(s)=\operatorname{Re}\int_0^\infty F(y)
        [\log(\sqrt N(s+iy))]^2\,dy,\qquad s\ge0 .
\]
Both the kernel and its derivatives decay exponentially at infinity.
The Fricke relation \(F(1/(Ny))=Ny^2F(y)\) gives flatness to every
order at zero. These bounds justify differentiation from the right
at the endpoint and the exchanges of absolutely convergent integrals
below.
\paragraph{Mellin normalization.}
For
\[
 \Lambda_E(w)=
 \left(\frac{\sqrt N}{2\pi}\right)^w
 \Gamma(w)L(E,w),
\]
Mellin transformation gives
\[
 \Lambda_E(w)=N^{w/2}
       \int_0^\infty F(y)y^{w-1}\,dy.
\]
At \(s=0\),
\[
 \operatorname{Re}\log^2(i\sqrt N y)
 =\log^2(\sqrt N y)-\frac{\pi^2}{4}.
\]
Since
\(\int_0^\infty F(y)\,dy=L(E,1)/(2\pi)=0\),
differentiation of the Mellin integral and the central vanishing yield
\[
 \boxed{\displaystyle
 P(0)=\frac{\Lambda_E''(1)}{\sqrt N}
     =\frac{L''(E,1)}{2\pi}.}
\]


\paragraph{Uniform positivity away from zero.}
Define
\[
 A_2=\int_0^\infty y^2F(y)\,dy
     =\frac{2L(E,3)}{(2\pi)^3},
 \qquad
 A_{4,\mathrm{abs}}=\int_0^\infty y^4|F(y)|\,dy.
\]
For \(s\ge1\), put \(B=\log(s\sqrt N)\), \(t=y/s\).
The exact identity
\[
 \operatorname{Re}\log^2(\sqrt N(s+iy))
 =\left(B+\frac12\log(1+t^2)\right)^2
       -\arctan^2t
\]
has expansion \(B^2+(B-1)t^2+R_B(t)\), where
\[
 R_B(t)=B\{\log(1+t^2)-t^2\}
       +t^2-\arctan^2t
       +\frac14\log^2(1+t^2).
\]
For \(t\ge0\), elementary integration gives
\[
 0\le t^2-\log(1+t^2)\le\tfrac12t^4,
 \quad
 0\le t^2-\arctan^2t\le\tfrac23t^4,
 \quad
 \log^2(1+t^2)\le t^4.
\]
Since \(B>0\),
\[
 |R_B(t)|\le
 \left(\frac B2+\frac{11}{12}\right)t^4.
\]
The \(B^2\) term integrates to zero by the central vanishing. Therefore
\[
 \boxed{\displaystyle
 \left|P(s)-\frac{(B-1)A_2}{s^2}\right|
 \le\frac{(B/2+11/12)A_{4,\mathrm{abs}}}{s^4}.}
\]

Euler products give explicit bounds. The Euler factors and
\(|a_n|\le d(n)\sqrt n\) imply
\[
 A_2\ge\frac{2}{(2\pi)^3}
       \left(\frac{\zeta(5)}{\zeta(5/2)}\right)^2
       >0.0048,
\]
\[
 A_{4,\mathrm{abs}}
 \le\frac{24}{(2\pi)^5}\zeta(9/2)^2
       <0.0028.
\]
For the first inequality, at every good prime the Euler
factor is bounded below by
\((1+\ell^{-5/2})^{-2}\); the multiplicative bad-prime
factors satisfy the same lower bound. For the second,
integrate the absolutely majorized Fourier series termwise.
For example,
\[
 \zeta(5)>1.036,\quad
 \zeta(5/2)<1.342,\quad
 \zeta(9/2)<1.055
\]
suffice. These rational bounds follow by finite summation and
\[
 \sum_{n>J}n^{-\sigma}
 \le \frac{J^{1-\sigma}}{\sigma-1}
 \qquad(\sigma>1).
\]

Let \(B_1=\log\sqrt N\), so \(6.18<B_1<6.19\).
The preceding three estimates give
\[
 P(1)>
 5.18(0.0048)
 -\left(\frac{6.19}{2}+\frac{11}{12}\right)(0.0028)
 >0.0136.
\]
Moreover,
\[
 \frac{B/2+11/12}{B-1}
\]
is strictly decreasing for \(B>1\).
Since \(B\ge B_1\) and \(s^{-2}\le1\), the remainder bound proves
\[
 \boxed{P(s)>0\quad\text{for every real }s\ge1.}
\]
It also gives the asymptotic
\[
 P(s)=\frac{A_2}{s^2}
       \{\log(s\sqrt N)-1\}
       +O\left(\frac{\log s}{s^4}\right),
 \qquad s\longrightarrow+\infty.
\]
In particular \(P(n)>0\) for every positive integer \(n\),
but \(P(n)\to0\). These nonzero values cannot all belong
to one fixed discrete additive subgroup of \(\mathbb R\).

\paragraph{The endpoint derivative.}
Differentiating the defining integral from the right at zero gives
\[
 P'(0)=\pi\int_0^\infty\frac{F(y)}y\,dy.
\]
The positive functional equation
\(\Lambda_E(w)=\Lambda_E(2-w)\) now yields
\[
 \boxed{\displaystyle
 P'(0)=\pi\Lambda_E(0)
       =\frac{N}{4\pi}L(E,2)>0.}
\]
The Euler-product estimate
\[
 L(E,2)\ge
 \left(\frac{\zeta(3)}{\zeta(3/2)}\right)^2
\]
gives the conservative explicit bound \(P'(0)>3900\).
Thus \(P(0)=0\), if true, is perfectly compatible with
strict positivity immediately to its right. No continuity
or one-sided positivity argument proves \(P(0)\ne0\).

\paragraph{Exact inverse Laplace density.}
Define the cosine transform
\[
 B_E(x)=\int_0^\infty F(y)\cos(xy)\,dy.
\]
The Laplace identity, for \(\operatorname{Re}z>0\),
\[
 \int_0^\infty e^{-zx}(-\log x-\gamma)\,dx
       =\frac{\log z}{z},
\]
where \(\gamma\) denotes Euler's constant, gives
\[
 P'(s)=-2\int_0^\infty e^{-sx}
       [\log(x/\sqrt N)+\gamma]B_E(x)\,dx.
\]
Integrating from \(s\) to infinity, using the large-translate asymptotic, proves
\[
 \boxed{\displaystyle
 P(s)=\int_0^\infty e^{-sx}W(x)\,dx,\qquad
 W(x)=\frac{2}{x}
       [\log(x/\sqrt N)+\gamma]B_E(x).}
\]
This also holds at \(s=0\). Because
\(\int F=0\), Taylor expansion at \(x=0\) gives
\[
 B_E(x)=-\frac{A_2x^2}{2}+O(x^4),
 \qquad W(x)=O(x|\log x|).
\]
At infinity \(W\) decreases rapidly, so \(W\in L^1(0,\infty)\).
There is no missing boundary mass at \(x=0\).

\paragraph{Infinitely many sign changes.}
By the Fricke relation, the even extension
\(\widetilde F(y)=F(|y|)\), with
\(\widetilde F(0)=0\), is Schwartz and flat at the origin.
Its Fourier transform is \(2B_E(x)\).
Fourier inversion and differentiation therefore imply
\[
 \boxed{\displaystyle
 \int_0^\infty x^{2j}B_E(x)\,dx=0
 \quad(j=0,1,2,\ldots).}
\]
The function \(B_E\) is not identically zero, since
\(B_E''(0)=-A_2<0\).

If \(B_E\) had finitely many sign changes, a polynomial
\(Q(x^2)\), with one simple root at each sign-change
location, could be assigned a global sign so that
\(Q(x^2)B_E(x)\ge0\) everywhere and is positive on an
open interval. Its integral would be positive. Expanding
\(Q\) would contradict the even-moment identities. Hence \(B_E\) has infinitely
many sign changes, necessarily at unbounded \(x\).
The logarithmic factor in the density has only one positive zero,
so \(W\) also changes sign infinitely often.

By uniqueness of Laplace transforms \cite{Widder} of finite signed
measures and Bernstein's characterization of completely
monotone functions, \(P\) is not completely monotone on
any half-line \(s>s_0\ge0\). Indeed, complete monotonicity
there would give a nonnegative representing measure for
\(P(s_0+\cdot)\), whereas its unique representing measure
is \(e^{-s_0x}W(x)\,dx\), which changes sign.

\paragraph{Finite-dimensional translation obstruction.}
There cannot be a finite-dimensional complex vector
space \(V\), a fixed \(T\in\operatorname{End}(V)\),
\(v\in V\), and \(\ell\in V^\vee\) satisfying
\[
 P(n)=\ell(T^nv)\qquad(n\ge1).
\]
Such a representation imposes a linear recurrence by Cayley--Hamilton.
Its generating function \(H(z)=\sum_{n\ge1}P(n)z^n\) would be rational.
The large-translate estimate gives
\[
 \sum_{n\ge1}P(n)<\infty,\qquad
 \sum_{n\ge1}nP(n)=\infty.
\]
Consequently \(H(r)\) has a finite limit as \(r\uparrow1\),
while \(H'(r)\to+\infty\). A rational function with a finite
limit at \(1\) has a removable singularity there and a
finite derivative, a contradiction. Thus a fixed
finite-dimensional linear translate orbit cannot encode
these raw periods. This statement does not exclude a
different arithmetic extension or a nonhorizontal
motivic construction.

\paragraph{Signed binomial reconstruction and Hausdorff
differences.}
For \(M\ge1\), define
\[
 S_M=\sum_{n=1}^M(-1)^{n+1}\binom MnP(n).
\]
The Laplace representation gives the exact identities
\[
 S_M=\int_0^\infty
       [1-(1-e^{-x})^M]W(x)\,dx,
\]
\[
 P(0)-S_M=
       \int_0^\infty(1-e^{-x})^M W(x)\,dx.
\]
Dominated convergence, using \(W\in L^1\), yields
\[
 \boxed{\displaystyle
 P(0)=\lim_{M\to\infty}
       \sum_{n=1}^M(-1)^{n+1}\binom MnP(n).}
\]
Thus the positive-index periods determine \(P(0)\),
but only through a signed reconstruction whose error
has no prescribed sign.

With \(\Delta P(n)=P(n+1)-P(n)\), we also have
\[
 (-1)^k\Delta^kP(n)=
 \int_0^\infty e^{-nx}(1-e^{-x})^kW(x)\,dx.
\]
These expressions cannot all be nonnegative for
\(n\ge1,\ k\ge0\). Otherwise, pushing the finite signed
measure \(e^{-x}W(x)\,dx\) forward under \(t=e^{-x}\)
would give a measure with nonnegative integrals against
every \(t^j(1-t)^k\). Approximation by Bernstein
polynomials would force that measure to be nonnegative,
contrary to the sign changes of \(W\).

Nevertheless, for every \emph{fixed} \(k\), the small-\(x\) expansion and
Laplace asymptotics give
\[
 (-1)^k\Delta^kP(n)
 \sim\frac{A_2(k+1)!}{n^{k+2}}\log n>0
 \qquad(n\to\infty).
\]
Thus each fixed finite-difference order eventually appears
positive, while the complete Hausdorff inequalities fail.


\paragraph{A smooth model with zero endpoint.}
The constraints above admit a Fricke-symmetric example with a vanishing
second central moment. Take a nonnegative even \(C^\infty\) function
\(h\) of integral one supported in \([-1/100,1/100]\), put
\(a=8/3\), and define
\[
G(t)=27\{3h(t)-2h(t-a)-2h(t+a)
               +\tfrac12 h(t-2a)+\tfrac12 h(t+2a)\}.
\]
Its bilateral Laplace transform is
\(54\widehat h(u)(\cosh(au)-1)^2\), where
\(\widehat h(u)=\int h(t)e^{ut}dt\).
Set \(F_*(y)=e^{-t}G(t)\) for \(t=\log(\sqrt N y)\). Evenness gives
the same Fricke identity, and its completed Mellin transform
\(N^{(1+u)/2}\int_0^\infty F_*(y)y^u\,dy\) has exact order four.
The corresponding \(P_*(0)\) is zero. Its quadratic and absolute
fourth ordinary moments satisfy
\[
 A_{2,*}\ge \frac{54(\cosh(2a)-1)^2}{N^{3/2}}>0.0050,\qquad
 A_{4,\mathrm{abs},*}
 \le\frac{27e^{0.04}(3+4\cosh(4a)+\cosh(8a))}{N^{5/2}}
 <0.00098.
\]
These follow by integrating the five translates; the absolute bound
uses the triangle inequality. The remainder estimate for \(P\) thus
gives \(P_*(s)>0\) for \(s\ge1\). This example tests the analytic
conditions only. It makes no assertion about Hecke eigenforms.

\section{A closed quadratic-twist period}
Let \(\chi(n)=(\frac{n}{3})\) and \(\omega_f=2\pi i f_E(z)\,dz\).
The normalized period of the cycle from \(2/3\) to \(1/3\) is
\[
\int_{2/3}^{1/3}\omega_f=i\sqrt3\,L(E\otimes\chi,1)
                         =16i\Omega^-,
\]
where \(i\Omega^-\) is the primitive purely imaginary period of
\(dx/(2y+x)\). The difference of boundary values of
\(\sum a_n e^{2\pi inz}/n\) gives the quadratic Gauss sum, using Abel
limits. Both cusps represent \(0\) at level \(N\), since \(3\nmid N\).
Cremona's optimality and Manin entry is \(1,1\) for the displayed
curve \cite{Cremona}. Its positive discriminant gives a rectangular
period lattice, so the quotient by \(i\Omega^-\) is an integer.

The twist has conductor \(9N\) and root number \(+1\), since
\(\chi(-N)=+1\). Put \(c=2\pi/(3\sqrt N)\). The Fricke-split Mellin
formula \cite[Ch.~II, Prop.~2.11.2]{Book} gives
\[
 L(E\otimes\chi,1)=2\sum_{n\ge1}\frac{a_n\chi(n)}n e^{-cn}.
\]
Use \(a_2=-1\); at the good odd primes \(p\le5000\) use
\(a_p=-\sum_{x\bmod p}(\frac{4x^3-3x^2-316x+1156}{p})\).
Extend by the Hecke recurrences and multiplicativity. The sum up to 5000
lies in \((7.395855,7.395858)\); the absolute tail is at most
\(4e^{-5001c}/(1-e^{-c})<0.00000038\) by \(|a_n|\le2n\).

Here are integer data sufficient to repeat the enclosure. Use
\(B=10^{25}\) and start with integer endpoints \(B,B\). Generate
successive bounds on \(Be^{-cn}\) by multiplying the previous endpoints by
\[
9956838327389655378430842/B,\quad
9956838327392406092287538/B,
\]
respectively, rounding out to integers. These bound \(e^{-c}\) by
the degree-seven and degree-six Taylor polynomials, using
\(3.1415926535<\pi<3.1415926537\) and bracketing \(\sqrt N\) to
width \(10^{-10}\). Multiply by \(2a_n\chi(n)/(nB)\), reverse the
endpoints when negative, and sum the intervals.

For roots \(e_1>e_2>e_3\) of \(4x^3-3x^2-316x+1156\), the
elliptic period formula gives
\[
 \Omega^-=\frac{\pi}
 {\operatorname{AGM}(\sqrt{e_1-e_3},\sqrt{e_2-e_3})}.
\]
Bisect the intervals \((5.74,5.75),(5.02,5.03),(-10.02,-10.01)\)
to width \(10^{-10}\). Five arithmetic and geometric mean steps,
bounding the limit between the final two means, give
\(0.80062<\Omega^-<0.80063\).
The resulting integer lies strictly between \(15.99\) and \(16.01\).

Quadratic base change to \(\mathbb Q(\sqrt{-3})\) therefore multiplies
the second central derivative by \(16\Omega^-/\sqrt3\), by the product
factorization and the two initial vanishings of \(L(E,s)\).

\section{A four-zero contour calculation}
For an even holomorphic function \(H\) assume that inside
\(|u|=R\) its zeros, with multiplicity, are \(0,0,\alpha,-\alpha\)
and that it has no boundary zero. Write \(\beta=\alpha^2\), allowing
\(\alpha=0\), and define the positively oriented integrals
\[
 M=\frac1{2\pi i}\oint u^2H'(u)/H(u)\,du,\qquad
 J_j=\frac1{2\pi i}\oint u^{2j+1}/H(u)\,du\quad(j=0,1).
\]
Then
\[
 M=2\beta,\qquad H''(0)=-\frac{M}{J_1-\beta J_0}.
\]
In fact \(H(u)=u^2(u^2-\beta)h(u^2)\) with \(h\) zero-free.
For \(k=1/h\), residues give \(J_1=k(\beta)\) and
\(J_0=(k(\beta)-k(0))/\beta\), interpreted as \(k'(0)\) if
\(\beta=0\). The denominator of the displayed expression is \(k(0)\).
For \(\alpha\ne0\), every bounded holomorphic weight \(\varphi\)
in the disc satisfying \(\varphi(0)=1\), \(\varphi(\pm\alpha)=0\)
obeys \(\|\varphi\|_\infty\ge R^2/|\beta|\).
Division by the two Blaschke factors proves the inequality, since their
product at zero has modulus \(|\beta|/R^2\). Normalizing that
product to one at zero attains equality. The stability cost is
therefore exact. Application to a particular disc for
\(\Lambda_E(1+u)\) requires a certified zero count on that disc.

\paragraph{Arithmetic lower bounds require an archimedean comparison.}
Finite congruence evidence for a regulator in another completion cannot
supply such a lower bound on its own. For instance
\(U(X)=X^4-X^3-X^2-X+1\) is irreducible modulo \(2\) (the fifth
cyclotomic polynomial there). Writing \(z=X+1/X\) gives
\(U(X)/X^2=z^2-z-3\). The smaller root \(z=(1-\sqrt{13})/2\)
is in \((-2,2)\), so one complex embedding of a root \(\alpha\)
satisfies \(\log|\alpha|=0\).
At \(17\), the simple root \(10\bmod17\) lifts to \(61\bmod289\),
and \(61^{16}\equiv1+6\cdot17\bmod289\).
In that embedding \(\log_{17}(\alpha^{16})/17\equiv6\bmod17\).
This verifies the directional limitation of a single congruence,
without addressing possible special comparisons for the elliptic
curve.

\begin{thebibliography}{9}
\bibitem{Cremona} J. E. Cremona, \emph{Elliptic curve data},
\href{https://github.com/JohnCremona/ecdata/blob/master/allcurves/allcurves.230000-239999}{curve and rank data}
and \href{https://github.com/JohnCremona/ecdata/blob/master/opt_man/opt_man.230000-239999}{optimality and Manin constants}, rows for conductor 234446.
\bibitem{GZK} A. Wiles, \emph{The Birch and Swinnerton-Dyer conjecture},
\href{https://www.claymath.org/wp-content/uploads/2022/05/birchswin.pdf}{Clay problem description}. See the results of Gross--Zagier and Kolyvagin on analytic ranks zero and one.
\bibitem{HLP} E. W. Howe, F. Lepr\'evost and B. Poonen,
\emph{Large torsion subgroups of split Jacobians of curves of genus two or three}, Forum Math. \textbf{12} (2000), 315--364,
\href{https://math.mit.edu/~poonen/papers/large.pdf}{Prop.~3}.
\bibitem{Book} J. E. Cremona, \emph{Algorithms for Modular Elliptic Curves},
\href{https://johncremona.github.io/book/fulltext/chapter2.pdf}{Chapter~II}.
\bibitem{Widder} D. V. Widder, \emph{The Laplace Transform}, Princeton University Press, 1941.
\end{thebibliography}
\end{document}
